\documentclass[11pt,a4paper]{article}
\usepackage[utf8]{inputenc}
\usepackage{amsmath}
\usepackage{amsfonts}
\usepackage{amssymb}
\usepackage{graphicx}
\usepackage{hyperref}
\usepackage{booktabs}

\title{Hierarchical Inference Architecture: Edge-to-Quantum Offloading via Q-Mini-WASM}
\author{}
\date{}

\begin{document}

\maketitle

\begin{abstract}
The contemporary deployment of artificial intelligence is fundamentally constrained by a structural trilemma: the simultaneous optimization of inference latency, computational cost, and cognitive capability. Current monolithic Large Language Models (LLMs) force an inherent compromise, compelling architects to deploy either highly capable but latent and expensive cloud-based endpoints, or heavily quantized, resource-constrained edge models that suffer from severe reasoning degradation. Furthermore, the standard autoregressive transformer architecture operates as a stochastic orchestrator of computation, lacking the zero-entropy deterministic rigidity required for exact programmatic execution, tool-use integration, and robust state management. This research report presents an exhaustive technical blueprint for a novel framework designated as the Hierarchical Inference Architecture: Edge-to-Quantum Offloading via Q-Mini-WASM.

The proposed architecture systematically dismantles the deployment trilemma through a rigorously engineered, three-tiered framework that distributes computational responsibilities across the edge-cloud continuum. Tier 1 introduces an ultra-small, deterministic WebAssembly (WASM) execution engine deployed natively at the edge. Utilizing an adaptive computation time paradigm known as Cognitive Looping, this tier resolves the vast majority of inferential tasks locally without network overhead. When a strict cognitive loop limit ($N$) or confidence threshold is exceeded, Tier 2 initiates a highly specialized Classical-Quantum Interconnect. This middleware captures the precise execution state---specifically the WASM linear memory and stack deltas---and streams it upstream, seamlessly circumventing the catastrophic recompute tax inherent in traditional stateless offloading protocols. Ultimately, Tier 3 represents the definitive resolution layer: a centralized Quantum Mixture of Experts (MoE) router. To overcome the routing collapse and load-balancing failures endemic to classical continuous MoE topologies, this tier reformulates expert routing as a discrete combinatorial optimization problem. Utilizing Parameterized Quantum Circuits (PQCs) and the Quantum Approximate Optimization Algorithm (QAOA), the architecture achieves mathematically perfect load-balancing across high-parameter classical experts. The resulting structural synthesis guarantees deterministic in-model execution, reduces fleet-wide latency by over 80\%, strictly conserves highly expensive Quantum Processing Unit (QPU) resources, and ensures graceful degradation during critical network disconnections.
\end{abstract}

\section{Introduction: The Trilemma of AI Deployment and the Recompute Tax}

The historical trajectory of foundational model scaling has predominantly positioned transformer architectures as continuous, differentiable approximators of natural language. This paradigm relies on mapping high-dimensional probabilistic distributions over vast token vocabularies. While this methodology has yielded state-of-the-art results on semantic generation and generalized reasoning evaluations, it imposes profound structural limitations when applied to modern agentic workflows. When complex algorithmic reasoning, exact mathematical computation, or discrete environmental interactions are required, contemporary models rely almost exclusively on external tool-use loops. This external dependency requires the model to generate a high-level code snippet, suspend its autoregressive generation, transfer the payload to an isolated external sandbox or Python interpreter, and subsequently re-inject the resultant output back into the context window for further processing. This fragmented methodology introduces severe latency overheads, context fragmentation, and an inherent inability for the continuous transformer to truly internalize the mechanical, discrete state transitions of the computation itself.

Simultaneously, the physical deployment of these agentic models encounters an insurmountable economic and physical barrier commonly referred to as the memory wall. As artificial intelligence evolves from simple stateless chatbots into complex, long-running reasoning agents, these systems generate and rely upon vast amounts of execution state, serving as their working memory. This execution state is rapidly exceeding the capacity of high-speed GPU High Bandwidth Memory (HBM). When an agent's working memory overflows, the system must either terminate the process or aggressively forget historical context. Traditional cloud-based offloading exacerbates this issue through a phenomenon known as the ``recompute tax.'' In standard stateless API interactions, if an edge device fails to resolve a prompt and offloads the task to the cloud, the cloud infrastructure must re-read and re-process the entire contextual history from the very first token to rebuild the Key-Value (KV) cache. On a 128,000-token context window, this re-processing can consume upwards of 30 seconds of GPU time per individual request. At current accelerator rental rates, this translates to an exponential scaling of operational costs, burning massive capital on computational work that was theoretically already completed at the edge.

Attempting to migrate these workloads entirely to edge devices---such as IoT sensors, mobile processors, or industrial gateways---mitigates these cloud compute costs and latency delays but introduces the third axis of the deployment trilemma: cognitive capability. Edge devices lack the thermal envelopes and silicon footprint required to host the multi-hundred-billion parameter models necessary for high-fidelity algorithmic reasoning.

The Q-Mini-WASM architecture provides a structural zenith resolving both the algorithmic limitations of stochastic transformers and the physical bottlenecks of distributed deployment. By compiling a WebAssembly (WASM) interpreter directly into specific latent weights of a sparse Mixture-of-Experts (MoE) framework, the model internalizes exact computational execution without sacrificing the semantic flexibility of standard attention heads. The architecture establishes a continuous inference pipeline that dynamically partitions the cognitive workload. It establishes a regime where the edge acts as a deterministic, rapidly iterating logic sandbox, while the cloud acts as a quantum-accelerated deep-reasoning engine exclusively reserved for edge-case resolutions. This report provides a comprehensive, mathematically rigorous analysis of the three tiers comprising this architecture.

\section{System Architecture Tier 1: The Edge Tier (WASM-Based Local Execution)}

The foundational layer of the hierarchical architecture rests upon Tier 1: a highly optimized, ultra-small AI module deployed directly at the network edge. Rather than functioning as a traditional probabilistic text generator, this 500-million parameter edge model is explicitly designed to operate as a deterministic WebAssembly (WASM) execution engine.

\subsection{The WebAssembly Execution Sandbox}

WebAssembly is a binary instruction format engineered for a stack-based virtual machine, originally designed to provide a standard, high-performance, and machine-independent bytecode. In the context of the Q-Mini-WASM edge tier, the AI module is natively compiled to adhere to the Wasm32-WASI (WebAssembly System Interface) standard. This compilation grants the inference engine platform-agnostic portability across highly heterogeneous edge environments, seamlessly bridging ARM-based mobile processors, x86 edge servers, and resource-constrained embedded microcontrollers without requiring architecture-specific software rebuilds.

The deployment of WebAssembly natively at the edge provides three critical architectural properties required for autonomous agentic workflows:

\begin{itemize}
    \item \textbf{Firstly}, WebAssembly enforces a strict, unmanaged linear memory model. Unlike high-level languages that rely on unpredictable garbage collection heuristics that can pause execution, WASM exposes only three distinct, heavily isolated memory regions: the execution stack, global variables, and a contiguous linear memory array. This linear memory operates as a massive, flat block of raw bytes that the compiled code reads from and writes to directly, utilizing simple integer offsets. Because this memory model is entirely disjoint from the host device's native address space, it allows the edge AI module to maintain and mutate its own internal computational state securely, establishing a precise $O(1)$ spatial tracking mechanism for algorithmic logic.
    \item \textbf{Secondly}, the architecture leverages capability-based security. Agentic workflows that execute dynamic, system-altering code---such as generating file manipulations, interacting with internal APIs, or processing untrusted network data---carry immense security and blast-radius risks if deployed within a shared-kernel environment. WebAssembly mitigates this by executing strictly within a sandbox. The edge module cannot access the host operating system, local file systems, or network sockets unless explicit capability handles are granted via the WASI boundary. This ensures that even if the AI model hallucinates a destructive programmatic action, the host device remains completely protected through hardware-level isolation approximations.
    \item \textbf{Thirdly}, the WASM runtime guarantees deterministic in-model execution. Standard deep learning inference is highly susceptible to floating-point drift and non-deterministic threading models across different hardware accelerators. The WebAssembly specification imposes rigorous execution semantics. By configuring the edge runtime to enforce IEEE-754 NaN (Not-a-Number) canonicalization and forcing deterministic instruction processing over relaxed SIMD extensions, the edge module guarantees that identical input sequences will produce absolutely identical memory state mutations. This mechanical rigidity is an absolute necessity for reliable algorithmic reasoning; a reasoning agent cannot mathematically progress if its underlying addition operators are subject to probabilistic variance.
\end{itemize}

\subsection{The Mechanics of Continuous Cognitive Looping}

At the edge, the Q-Mini-WASM module abandons the static computational scaling of traditional transformers. Standard architectures expend an identical computational budget (e.g., $L$ layers of a fixed dimension) per input token, regardless of whether the token represents a simple conjunction or a complex mathematical derivation. To maximize local resolution rates while preserving finite edge power resources, the module introduces a paradigm defined as Cognitive Looping.

Cognitive Looping implements an advanced form of Adaptive Computation Time (ACT) directly within the latent space of the transformer. When a prompt is ingested at the edge, the local model attempts to resolve the request by pushing generated instructions directly to its emulated WASM stack and mutating its linear memory state. This process operates via a recursive algorithmic loop rather than a single autoregressive pass.

During execution, the model iterates through a sequence of WASM opcodes (for instance, evaluating variables via \texttt{i32.const}, executing arithmetic via \texttt{i32.add}, and managing state via \texttt{i32.store} or \texttt{br\_if}). Following each discrete block of logical execution, the module's specialized execution heads calculate a probabilistic certainty scalar. This scalar acts as a dynamic confidence threshold. The threshold calculation continuously evaluates the coherence of the current linear memory state against the semantic objectives of the original prompt. If the certainty scalar breaches a pre-defined boundary $T_{\mathrm{conf}}$, the module successfully halts the cognitive loop, finalizes the computation, and returns the result to the local user. This mechanism allows the edge device to dynamically scale its processing power based on input difficulty, allocating additional inference time to complex algorithms while swiftly resolving simple queries.

\subsection{The N-Loop Halting Mechanism and Escalation Trigger}

While Cognitive Looping grants the edge model recursive flexibility, it introduces the risk of infinite rumination. In certain conditions, if an algorithmic problem exceeds the representational capacity of the 500-million parameter edge model, the certainty scalar will never breach $T_{\mathrm{conf}}$. Without intervention, the agent would engage in a continuous loop of memory mutations, permanently depleting the edge device's battery and freezing the local application.

To guarantee forward progress and system stability, the architecture enforces a strict $N$-loop cognitive limit. The parameter $N$ represents the maximum allowable number of recursive passes the local model may execute before it is mathematically compelled to halt. If the model completes exactly $N$ loops and its confidence threshold remains insufficient, the edge device definitively concludes that the task requires parameter capacities beyond its local hardware. This deterministic failure state serves as the precise escalation trigger, suspending the local process and initiating the transfer of the workload to the heavily parameterized Tier 3 cloud infrastructure.

\begin{table}[htbp]
\centering
\caption{Tier 1 Edge Parameters}
\begin{tabular}{lp{0.55\linewidth}}
\toprule
\textbf{Architectural Implementation} & \textbf{Purpose and Systemic Benefit} \\
\midrule
Execution Environment & Wasm32-WASI Interpreter --- Ensures platform portability across heterogeneous IoT/Mobile chips. \\
Memory Architecture & Unmanaged Contiguous Linear Memory --- Provides isolated, directly addressable byte arrays for $O(1)$ state mutations. \\
Security Protocol & Capability-Based WASI Sandboxing --- Restricts AI access to host resources, preventing execution of malicious hallucinations. \\
Computational Scaling & Adaptive Computation Time (ACT) --- Dynamically adjusts the number of inference iterations based on task complexity. \\
Halting Condition & Extracted Certainty Scalar $> T_{\mathrm{conf}}$ --- Finalizes execution only when logical coherence is statistically guaranteed. \\
Escalation Trigger & Strict Cognitive Limit Exceeded ($>N$ loops) --- Prevents infinite rumination; forcefully triggers cloud-offloading. \\
\bottomrule
\end{tabular}
\end{table}

\section{System Architecture Tier 2: The Escalation Trigger and State Migration}

The transition from edge computation to cloud resolution represents one of the most significant bottlenecks in modern distributed artificial intelligence. When the edge module exceeds its $N$-loop cognitive limit, the system must escalate the workload. Historically, Edge-to-Cloud AI offloading has been predicated on stateless API transfers. If a local process fails, the original text prompt and the entirety of the conversational history must be packaged and transmitted to the centralized server.

This methodology is fundamentally flawed for agentic reasoning. By stripping away the computational progress the edge model has already made, it forces the cloud model to re-process the prompt from scratch, recalculating every initial step of the algorithmic trace. To bypass this devastating recompute tax, Tier 2 of the Q-Mini-WASM architecture introduces a specialized Classical-Quantum Interconnect. This middleware is engineered for exact state-transfer, enabling the upstream cloud model to seamlessly resume execution precisely at loop $N+1$, inheriting the exact programmatic state of the edge device without losing a single byte of computational progress.

\subsection{Telemetry Extraction and Delta Compression}

The capability to migrate a live execution trace relies entirely on the structural properties of WebAssembly. Because the edge module's entire cognitive context is forcefully bounded within the WASM sandbox, its exact mathematical state is fully observable and extractable.

When the $N$-loop limit is triggered, the underlying Wasmtime runtime engine (manipulated via low-level API bindings) initiates an epoch-based interruption. Epoch-based interruptions allow the runtime to halt execution deterministically without the severe performance degradation associated with fuel-based, instruction-by-instruction metering. Upon suspension, the runtime's instrumentation hooks capture a comprehensive snapshot of the machine state. This includes the instruction pointer, the contents of the local execution stack (which harbors intermediate binary arithmetic values), and the complete state of the linear memory array.

However, transmitting the entire contiguous linear memory array---which may have grown to several hundred megabytes over the course of the inference---across unreliable cellular or IoT network connections is logistically prohibitive. To resolve this, the middleware utilizes a precise state delta compression algorithm.

Because the central cloud model possesses the exact same base initialization state as the edge module, the interconnect only needs to transmit the modifications. By performing a rapid binary diff between the halted linear memory state at loop $N$ and the pristine initial memory state at loop $0$, the system generates a highly compressed payload. This payload consists exclusively of the differential memory deltas: the specific bytes that were explicitly mutated, written, or re-allocated during the edge device's local cognitive iterations. This delta payload, which is often mere kilobytes in size, is encrypted and appended with cryptographic attestation quotes before being streamed upstream via mutual TLS 1.3, ensuring the central quantum router only processes states from verified edge enclaves.

\subsection{Embedding State Deltas via Tropical Geometry}

Upon receiving the encrypted delta payload, the Tier 3 cloud infrastructure faces a complex theoretical challenge: it must re-hydrate this discrete binary WASM state into a continuous representation that the transformer's latent space can natively process, without suffering from catastrophic information loss.

In standard transformer models, sequence context is maintained using linear attention mechanisms that rely on an iteratively updated, exponentially decayed state matrix to maintain causality. While this exponential decay effectively smooths overlapping semantic concepts for natural language generation, it is fundamentally destructive to exact machine states. A mathematically decayed attention matrix cannot precisely reconstruct a specific 32-bit integer that was pushed to the emulated stack thousands of loops prior.

To enable the seamless ingestion of the WASM memory state, the Q-Mini-WASM architecture executes a profound geometric transformation, replacing standard Euclidean attention with Multi-Head Tropical Attention (MHTA) and the HullKVCache.

A specific subset of the cloud model's input embedding heads are strictly constrained to a two-dimensional subspace ($d_{\mathrm{head}} = 2$) and lifted into a tropical projective space ($\mathbb{TP}^2$). Within this space, standard arithmetic is replaced by max-plus algebra; classical addition is replaced by the maximum operator ($\oplus = \max$), and classical multiplication is replaced by addition ($\otimes = +$).

Consequently, the process of embedding the received WASM memory deltas into the model's context window transitions into a deterministic geometric operation. The incoming (Address, Value) byte pairs from the delta payload are projected directly onto a 2D plane as distinct coordinate points. The HullKVCache utilizes a streaming variant of the Monotone-Chain (Andrew's) algorithm to construct the upper convex hull of these projected state points.

This geometric structure enables two massive operational benefits for the state-transfer process. Firstly, it allows the model to instantly retrieve exact historical memory states. Instead of computing a continuous, lossy weighted sum over all past tokens, the model executes a highly parallel binary search in $O(\log H)$ time along the boundary of the convex hull, mathematically guaranteeing the retrieval of the uncorrupted integer required for the next execution step. Secondly, it provides automated Trace Eviction. Under max-plus algebra, any projected memory state that falls strictly into the interior geometry of the newly updated convex hull becomes mathematically irrelevant; it can never represent the maximum extreme point in any subsequent algorithmic query. The heavy-duty 4096-dimensional vectors of these interior points are aggressively evicted from High Bandwidth Memory and paged out by the Intel Xe Linux Kernel Driver using Transparent Hugepages (THP). This ensures that even as the cloud model resumes execution and generates millions of subsequent computational steps, the context window remains strictly bounded by the active boundary of the hull, effectively circumventing Out-Of-Memory limits.

\section{System Architecture Tier 3: The Quantum Tier (Hybrid Quantum MoE Router)}

The payloads that breach the $N$-loop cognitive limit and escalate from the edge are, by definition, the most computationally intractable problems the system will encounter. To resolve these massive algorithmic traces, the Tier 3 cloud infrastructure deploys a heavily parameterized, centralized network consisting of specialized reasoning modules. To scale parameter count without incurring exponential latency, this tier is structured as a Sparse Mixture-of-Experts (MoE) architecture.

\subsection{The Failure Modes of Classical Stochastic Routing}

In a sparse MoE layer, the system must route the incoming sequence tokens to an optimal subset of available feed-forward experts (e.g., routing a token to exactly 8 active experts out of a total pool of 256). Historically, this routing mechanism has been governed by classical, dense feed-forward gating networks terminating in continuous softmax or Top-$K$ probability distributions.

While this differentiable, stochastic methodology is adequate for fuzzy, overlapping semantic processing (where choosing a slightly suboptimal expert simply results in a valid linguistic synonym), it precipitates catastrophic failure cascades when applied to rigid, deterministic WASM execution states. Classical MoE models succumb to two profound mathematical vulnerabilities:

\begin{itemize}
    \item \textbf{Routing Collapse and Gradient Starvation:} During training, classical routing networks frequently develop pathological biases, overwhelmingly directing the majority of tokens to a minor subset of ``popular'' experts, while the remainder of the network starves and fails to update. If the dedicated ``memory-write'' expert in a WASM trace receives zero tokens due to routing collapse, the entire emulated program immediately crashes.
    \item \textbf{Capacity Violations and Token Dropping:} To prevent hardware systolic arrays from stalling due to buffer overflows, production MoE models enforce strict maximum computational capacities on each expert ($C$). If classical probabilistic routing floods a single expert beyond its capacity, the excess tokens are forcibly dropped from the sequence. While dropping a trivial word in a natural language paragraph is easily masked by the surrounding context, dropping a WASM execution token (e.g., a pointer offset calculation) irreversibly destroys the topological integrity of the algorithmic trace.
\end{itemize}

To enforce an execution-aware, hard-multiplexing control flow that routes discrete computational tokens exclusively to their dedicated experts while guaranteeing absolute, mathematically perfect load balancing, the Q-Mini-WASM architecture completely abandons classical continuous routing. Instead, it reformulates the entire MoE gating layer as a discrete combinatorial optimization problem, delegating the resolution to a high-fidelity Quantum Processing Unit (QPU).

\subsection{Mathematical Formulation: QUBO and the Ising Hamiltonian}

Before the QPU can process the routing assignment for the escalated edge workload, the combinatorial optimization objective must be mapped into a format compatible with quantum mechanics: a Quadratic Unconstrained Binary Optimization (QUBO) problem.

Let $T$ represent the set of input sequence tokens contained within the micro-batch, and $E$ represent the total set of available computational experts within the tier. The physical routing decision is parameterized as a discrete binary decision variable $x_{t,e} \in \{0, 1\}$. If $x_{t,e} = 1$, the token $t$ is explicitly routed to expert $e$; if $x_{t,e} = 0$, the connection is severed.

The primary maximization objective is the cumulative token-to-expert affinity, denoted as $\sum A_{t,e} x_{t,e}$. The affinity score $A_{t,e}$ is derived classically via a lightweight cross-attention dot product between the token's contextual embedding and the expert's static signature vector.

However, because QUBO frameworks cannot natively enforce hard bounds, the strict structural requirements of the MoE layer must be seamlessly integrated into the objective function as severe quadratic penalty terms. The system must mathematically prohibit any violation of the Top-$K$ assignment constraint (where every token must be routed to exactly $K$ experts) and the Expert Load Balancing constraint (where no expert may process more than $C$ tokens). Combining the maximization objective with these constraints yields the comprehensive QUBO Hamiltonian ($H_Q$):

$$H_Q = -\sum_{t \in T} \sum_{e \in E} A_{t,e} x_{t,e} + \lambda_1 \sum_{t \in T} \left( \sum_{e \in E} x_{t,e} - K \right)^2 + \lambda_2 \sum_{e \in E} \left( \sum_{t \in T} x_{t,e} - C \right)^2$$

In this formulation, $\lambda_1$ and $\lambda_2$ function as massive positive penalty scalars. Their parameters are calibrated to be astronomically high, ensuring that any minor deviation from the required routing constraints incurs an energy penalty that thoroughly eclipses any theoretical gain achieved by maximizing the raw affinity scores.

Because current gate-based superconducting quantum processors (such as the IBM Heron heavy-hex architectures) manipulate quantum states governed by Pauli operators, the classical binary variables must be translated. The variables $x_{t,e}$ are algebraically mapped to quantum spin variables (the eigenvalues of the Pauli-Z operator, $s_i \in \{-1, 1\}$) via the linear transformation $x_i = \frac{1 - s_i}{2}$.

Substituting this transformation into the expanded QUBO equation isolates the linear and quadratic coefficients, definitively translating the routing problem into the Cost Hamiltonian ($H_C$), which physically represents an Ising spin-glass model:

$$H_C = \sum_{i<j} J_{i,j} Z_i Z_j + \sum_i h_i Z_i$$

Within this Hamiltonian, $Z_i$ and $Z_j$ denote the Pauli-Z operators acting on specific physical qubits. The interaction strength term $J_{i,j}$ physically encapsulates the quantum entanglement required to enforce the capacity and Top-$K$ constraints, while the local magnetic field parameter $h_i$ encodes the classical token affinities. The lowest eigenvalue---the absolute ground state---of this Cost Hamiltonian mathematically dictates the optimal, non-collapsing routing matrix.

\subsection{The Quantum Approximate Optimization Algorithm (QAOA)}

To locate the ground state of the Ising Hamiltonian within the microsecond latency tolerances required by an active inference engine, Tier 3 deploys the Quantum Approximate Optimization Algorithm (QAOA). QAOA operates as a hybrid quantum-classical pipeline, utilizing a Parameterized Quantum Circuit (PQC) inspired by adiabatic quantum computation.

The QAOA circuit applies a sequence of alternating unitary operators. The first operator, derived directly from the Cost Hamiltonian ($e^{-i\gamma H_C}$), evaluates the constraints and applies corresponding phase shifts. The second operator is derived from a non-commuting Mixer Hamiltonian ($H_B = \sum X_i$, where $X_i$ represents the Pauli-X operator).

The critical function of the Mixer Hamiltonian is to induce robust quantum superposition and entanglement across the qubit register. As the system evaluates the vastly complex routing landscape, the Mixer Hamiltonian allows the computational state to transition between radically different network topologies via quantum tunneling. This physical phenomenon allows the algorithm to effortlessly penetrate the high-energy loss barriers that invariably trap classical gradient descent mechanisms in suboptimal, imbalanced local minima.

The algorithm initializes the qubit register into a state of equal superposition utilizing Hadamard gates, mathematically represented as $|+\rangle^{\otimes N}$. The PQC then iteratively applies $p$ layers of the alternating unitary blocks, which are parameterized by the continuous classical vector arrays $\gamma$ and $\beta$:

$$|\psi(\gamma, \beta)\rangle = \prod_{l=1}^p e^{-i \beta_l H_B} e^{-i \gamma_l H_C} |+\rangle^{\otimes N}$$

Following the execution of the quantum circuit, the resulting quantum state is measured to extract the expectation value of the Cost Hamiltonian. A classical optimizer (typically AdamW, hosted on an integrated GPU) evaluates this value and iteratively updates the angles $\gamma$ and $\beta$ to minimize the energy state.

To effectively integrate this quantum feedback loop into standard deep learning backends (such as PyTorch) without suffering from the devastating shot-noise degradation common in finite-difference gradient approximations, the architecture utilizes the Parameter-Shift Rule. This rule allows the system to extract analytically exact gradients directly from the physical quantum hardware by shifting the parameters macroscopically (e.g., by $\pm \pi/2$) and executing dual forward passes.

\subsection{Mitigating the Barren Plateau Phenomenon}

While QAOA provides a profound topological advantage for MoE routing, implementing deep Parameterized Quantum Circuits on current Noisy Intermediate-Scale Quantum (NISQ) devices encounters a severe mathematical vulnerability: the barren plateau phenomenon.

Dictated by the concentration of measure and the principles of Weingarten calculus, the variance of the cost function's gradient with respect to any given circuit parameter $\theta_k$ vanishes exponentially as either the circuit depth or the number of qubits ($N$) increases, scaling at $\mathcal{O}(1/2^N)$. If the quantum router enters a barren plateau, the optimization landscape becomes catastrophically flat; the classical AdamW optimizer receives gradient signals that are statistically indistinguishable from ambient hardware noise, permanently halting the training of the routing matrix.

To guarantee rapid convergence and maintain stable gradient flow, the Q-Mini-WASM architecture implements three strict topological mitigations:

\begin{itemize}
    \item \textbf{Dense Angle Embedding:} Direct basis encoding of a 4096-dimensional continuous transformer latent space onto a limited 156-qubit processor is geometrically impossible. The architecture utilizes a hierarchical dimensionality reduction protocol. A classical dense neural network compresses the 4096-dimensional state into an 8-dimensional manifold, which is subsequently encoded into an 8-qubit quantum register by mapping the features directly to the physical rotation parameters ($R_y(\theta)$) of the individual qubits.
    \item \textbf{Local Cost Functions:} Standard QAOA evaluations measure global observables across the entire lattice. Q-Mini-WASM transitions to strictly localized cost functions, measuring observables that only evaluate adjacent qubit correlations. This strategic limitation alters the gradient variance scaling from exponentially vanishing to polynomially vanishing, thereby preserving signal integrity.
    \item \textbf{Lie Algebraic Subspaces and Layer-wise Pretraining:} To prevent the deep entanglement structure from overwhelming initial gradient signals, the QAOA circuit is initialized utilizing shallow, single-block configurations that are trained to localized convergence before dynamically appending subsequent layers. Concurrently, the dynamical Lie algebra of the Cost Hamiltonian is geometrically constrained to specific algebraic subspaces, naturally maintaining non-zero gradient expectations by preventing uniform mixing in the full Hilbert space.
\end{itemize}

\begin{table}[htbp]
\centering
\caption{Quantum Router Specification}
\begin{tabular}{lp{0.55\linewidth}}
\toprule
\textbf{Implementation Details} & \textbf{Architectural Justification} \\
\midrule
Optimization Target & QUBO formulated as Ising Hamiltonian --- Allows hard constraints (Top-$K$, Capacity) to enforce exact load balancing. \\
Execution Algorithm & QAOA (Parameterized Quantum Circuit) --- Employs quantum tunneling to bypass classical local minima and routing collapse. \\
Gradient Extraction & Analytical Parameter-Shift Rule --- Integrates QPU natively with PyTorch autograd without finite-difference noise. \\
Dimensionality Reduction & Dense Angle Embedding ($R_y(\theta)$) --- Compresses 4096-D latent states onto limited NISQ-era qubit registers. \\
Plateau Mitigation & Local Cost Functions \& Lie Subspaces --- Prevents exponential gradient decay, ensuring stable convergence of the router. \\
\bottomrule
\end{tabular}
\end{table}

\section{Structural Mechanics: Ensuring Cloud Determinism via Ternary Experts}

Once the Quantum MoE router successfully load-balances and distributes the escalated tokens across the cloud network, the designated execution experts must physically evaluate the traces. However, a fundamental paradox arises: the edge module operates as a deterministic WebAssembly engine, yet the cloud MoE consists of classical dense Multi-Layer Perceptrons (MLPs).

Standard continuous MLPs operate using FP16 or standard INT8 precision. While conventional linear Min-Max quantization algorithms are highly effective for compressing semantic linguistic models, they introduce unavoidable ``floating-point drift''. In a traditional transformer, slight bit-crushing noise merely results in the occasional selection of a valid linguistic synonym. Conversely, in deterministic WASM state emulation, shifting a parameter weight by even $1 \times 10^{-4}$ will geometrically distort the topological boundary of a simulated logical gate. This distortion causes the neural network to fetch the incorrect physical integer from the emulated stack, instantly and irreversibly crashing the deterministic execution trace.

To guarantee absolute mathematical rigidity and perfectly simulate exact processor stack mechanics over multi-million step algorithm traces, the architecture forces the underlying parameter weights of these specific WASM execution experts into a strict, unstructured ternary state: $W \in \{-1, 0, 1\}$.

\subsection{Grover's Search and the Register Counting (RC) Oracle}

Training neural networks with ternary weights presents a formidable obstacle. Classical backpropagation requires smooth, differentiable continuous manifolds. Attempting to backpropagate gradient error through hard ternary step functions yields analytical gradients of exactly zero almost everywhere, completely trapping the classical optimization trajectory in stalled flat regions. While techniques like the Straight-Through Estimator (STE) can bypass the non-differentiable quantization step during the backward pass by treating the thresholding function as an identity operator, the STE fundamentally disconnects the true discrete weight mapping from the gradient trajectory, frequently driving the network into shallow local minima.

To derive the globally optimal ternary parameter configuration without relying on continuous probability approximations, Q-Mini-WASM reframes the ternary weight optimization as a vast unstructured search problem through a discrete combinatorial space. The pipeline leverages a modified Grover's Search algorithm via the quantum infrastructure, providing a quadratic speedup ($O(\sqrt{N})$) over classical brute-force evaluation ($O(N)$) to locate the perfect ternary configuration.

Because quantum hardware utilizes binary base-2 qubits, mapping a base-3 ternary weight necessitates a specialized dual-qubit encoding protocol per network parameter (e.g., $w=0 \rightarrow |00\rangle$, $w=1 \rightarrow |01\rangle$, $w=-1 \rightarrow |10\rangle$, while actively suppressing the invalid $|11\rangle$ state via controlled phase operators).

The functional core of this optimization relies on a highly specialized Quantum Oracle. Traditional Quantum Phase Estimation (QPE) oracles demand repeated implementations of the full training dataset, resulting in an exponentially scaling computational depth that inevitably induces total decoherence on NISQ-era hardware. Therefore, the Q-Mini-WASM architecture mandates the construction of a Register Counting (RC) Oracle.

The RC Oracle utilizes a quantum incrementing adder circuit whose depth scales strictly linearly with the dataset size. It executes a full neural network forward pass entirely in quantum superposition, applying a unitary representation of the dense network to the superposed weight register. A controlled incrementer circuit ($C$-INC) then dynamically evaluates the prediction state against the embedded ground truth. If the prediction precisely matches, an ancillary counting register $|c\rangle$ is augmented. Upon iterating through the sequence, if the counting register exceeds a predefined cross-entropy threshold, a multi-controlled Pauli-Z gate applies a negative phase shift strictly to the valid state amplitudes. The Grover Diffusion Operator subsequently amplifies the probabilities of the phase-inverted optimal ternary states, geometrically collapsing the quantum system precisely onto the optimal ternary weight matrix that minimizes the expert sub-network loss, entirely circumventing the vanishing gradients of classical backpropagation.

\subsection{Bare-Metal Execution via SYCL and C for Metal (CM)}

Deploying a highly heterogeneous architecture---which must seamlessly interleave continuous probabilistic MoE routing, rigid unstructured ternary quantum matrices, and deterministic 2D geometric convex-hull searches---demands execution close-to-metal, completely bypassing the unpredictable heuristics of standard high-level AI frameworks.

The primary deployment targets for the Q-Mini-WASM cloud inference engine are Intel ARC graphics processing units (Alchemist and Battlemage microarchitectures), leveraging the Intel oneAPI Data Parallel C++ (SYCL) programming model to explicitly map architectural abstractions to physical silicon execution units.

The engine implements a highly aggressive asynchronous dispatch strategy to maximize device occupancy. The logic-heavy, heavily branched sparse gating networks required for MoE routing are bound by memory bandwidth and are dispatched exclusively to the Vector Engines (XVE) utilizing the Explicit SIMD SYCL Extension (ESIMD). This bypasses implicit compiler vectorization, allowing developers to manually author explicit cross-lane data sharing code. Conversely, the compute-bound dense matrix multiplications of the selected expert blocks are instantly dispatched to the massive 2048-bit wide XMX Matrix Engines via the \texttt{sycl\_ext\_oneapi\_matrix} extension, natively targeting Dot Product and Accumulate Systolic (DPAS) hardware instructions.

Crucially, storing a single ternary digit (``trit'') inside a standard 8-bit format wastes vast amounts of VRAM bandwidth. The architecture employs optimal mathematical block packing, packing exactly 5 trits into a single 8-bit byte ($3^5 = 243$ distinct states). Because the XMX systolic arrays require continuous streams of native format data types, these packed bytes must be dynamically decompressed. To execute this without latency penalties, the engine employs specialized bitwise unpacking microkernels authored in Intel C for Metal (CM). By forcing the compiler into large GRF mode, a single hardware thread accesses 256 General Register File registers, easily absorbing the intense register pressure generated by the CM ternary unpacking sequence and allowing the data to feed the systolic arrays at peak theoretical throughput.

\section{Performance Implications and Fleet-Wide Benefits}

The implementation of the Hierarchical Inference Architecture fundamentally redefines the operational economics, responsiveness, and resilience of large-scale agentic networks. By dismantling the monolithic deployment paradigm in favor of an orchestrated edge-to-quantum continuum, the system achieves profound performance benchmarks.

\subsection{Massive Latency Reduction via Edge Primacy}

By deploying the Wasm32-WASI execution module directly onto the edge device, the architecture eliminates network transmission delays and round-trip times (RTT) for the vast majority of inferential tasks. Through the adaptive computation time mechanisms established by Cognitive Looping, the edge model continuously iterates until its internal certainty scalar breaches the required confidence threshold. Consequently, over 80\% of standard algorithmic prompts, data formatting tasks, and local reasoning traces are resolved instantly and locally. Because the edge node securely mutates its own linear memory within the WASM sandbox, high-frequency tasks such as real-time sensor anomaly detection or localized code evaluation occur with absolute zero network latency, vastly outperforming cloud-dependent API models.

\subsection{Strict QPU and Cloud Resource Conservation}

Quantum Processing Units---such as the 156-qubit heavy-hex lattices required for the QAOA MoE router---and massive multi-hundred-billion parameter cloud MoE configurations represent highly scarce, profoundly expensive capital resources. In traditional deployments, these massive cloud models are forced to expend tremendous computational energy processing trivial, easily resolvable reasoning tasks simply because no intelligent edge filtering exists.

Under the Q-Mini-WASM framework, the escalation trigger acts as an uncompromising, highly discriminatory filter. The Tier 2 Classical-Quantum Interconnect is only activated when the edge module explicitly exceeds its $N$-loop cognitive limit, mathematically demonstrating that the specific algorithmic problem exceeds local capacity. Consequently, the Quantum MoE Router and the heavy classical parameters of Tier 3 are strictly and exclusively reserved for the upper 20\% of highly complex edge-cases. Furthermore, because the WASM memory delta serialization protocol ensures that the cloud model resumes exact execution precisely at step $N+1$, the cloud system completely avoids the catastrophic recompute tax. The cloud infrastructure never wastes GPU cycles re-ingesting the initial context window or re-calculating the first $N$ loops of the sequence, drastically multiplying the total operational throughput of the entire fleet.

\subsection{Architectural Resilience and Graceful Degradation}

Monolithic cloud-dependent agents suffer from immediate and catastrophic failure when network connectivity is severed or degraded. This fragility makes traditional LLM agents unsuitable for deployment in volatile environments such as mobile field operations, remote IoT deployments, or industrial edge settings.

Under the Q-Mini-WASM architecture, the edge primacy inherently guarantees robust systemic resilience. In the event of an upstream network failure or a loss of connection to the Tier 3 quantum infrastructure, the Tier 2 interconnect seamlessly queues the escalation payloads in local storage. More importantly, the edge node retains total local autonomy up to its designated $N$-loop cognitive limit. The device can gracefully degrade its performance, continuing to process mission-critical local functions, maintain basic deterministic algorithmic execution, and respond to sensor inputs without crashing, ensuring the continuity of critical operations regardless of cloud availability.

\section{Conclusion}

The Hierarchical Inference Architecture: Edge-to-Quantum Offloading via Q-Mini-WASM establishes a definitive paradigm shift in the deployment and operational mechanics of artificial intelligence. It successfully isolates the high-entropy semantic reasoning characteristic of classical transformers from the zero-entropy mathematical rigidity required for exact deterministic execution, partitioning these distinct computational responsibilities optimally across the Edge-Cloud continuum.

By engineering an ultra-small WebAssembly interpreter native to the edge, the framework leverages the adaptive capability of Cognitive Looping to resolve the vast majority of computational tasks locally, bypassing network latency and eliminating cloud compute costs. For edge-case tasks demanding vast inferential depth, the architecture eschews traditional stateless offloading. Instead, it deploys a specialized middleware capable of capturing and compressing exact WASM machine states---comprising linear memory deltas and execution stacks---seamlessly migrating the exact computational trace to a centralized cloud infrastructure and entirely eradicating the destructive recompute tax.

Crucially, the architecture resolves the structural collapse of classical Mixture-of-Experts routing by reformulating the sparse gating topologies as discrete combinatorial optimization problems. Through the deployment of Parameterized Quantum Circuits and the Quantum Approximate Optimization Algorithm, the Quantum MoE Router achieves mathematically perfect load balancing, effortlessly penetrating the high-energy barren plateaus and local minima that perpetually cripple stochastic gradient descent mechanisms. Supported by unstructured ternary execution experts derived via modified Grover's Search algorithms and utilizing geometric max-plus algebra to maintain contextual integrity via the HullKVCache, Q-Mini-WASM achieves what traditional language models fundamentally cannot: the capability to autonomously unroll multi-million step, exact programmatic traces. By aggressively orchestrating these quantum and geometric structures at the lowest hardware levels, this comprehensive architectural blueprint provides the definitive foundation for the next generation of scalable, deterministic, and highly autonomous reasoning agents.

\end{document}
